% Auto-generated answer document. Original problem statements are preserved verbatim.
\documentclass[11pt]{article}
\usepackage{fontspec}
\usepackage{microtype}
\usepackage{mathtools,amssymb,amsfonts,amsthm,bm}
\usepackage{enumitem}
\usepackage{geometry,fancyhdr,setspace,xcolor}
\usepackage{fancyvrb}
\usepackage{pdfpages}
\usepackage[hidelinks]{hyperref}
\geometry{margin=1in,headheight=14pt}
\setstretch{1.15}
\setlength{\parindent}{0pt}
\setlength{\parskip}{0.55em plus 0.12em minus 0.08em}
\setlength{\abovedisplayskip}{0.85em plus 0.2em minus 0.1em}
\setlength{\belowdisplayskip}{0.85em plus 0.2em minus 0.1em}
\setlist{leftmargin=*,topsep=0.35em,itemsep=0.15em,parsep=0pt}
\pagestyle{fancy}
\fancyhf{}
\fancyhead[L]{\coursename}
\fancyhead[C]{\assignment}
\fancyhead[R]{\studentname}
\fancyfoot[C]{\thepage}
\newcommand{\studentname}{Keshav Ramamurthy}
\newcommand{\coursename}{}
\newcommand{\assignment}{}
\newcommand{\term}{Fall 2026}
\newcommand{\R}{\mathbb{R}}
\newcommand{\C}{\mathbb{C}}
\newcommand{\Q}{\mathbb{Q}}
\newcommand{\Z}{\mathbb{Z}}
\newcommand{\N}{\mathbb{N}}
\newcommand{\F}{\mathbb{F}}
\newcommand{\K}{\mathbb{K}}
\newcommand{\E}{\mathbb{E}}
\newcommand{\Prob}{\mathbb{P}}
\newcommand{\abs}[1]{\left\lvert#1\right\rvert}
\newcommand{\norm}[1]{\left\lVert#1\right\rVert}
\newcommand{\ip}[2]{\left\langle#1,#2\right\rangle}
\newcommand{\set}[1]{\left\{#1\right\}}
\newcommand{\given}{\,\middle\vert\,}
\newcommand{\dd}{\,\mathrm{d}}
\DeclareMathOperator{\rank}{rank}
\DeclareMathOperator{\nullity}{nullity}
\DeclareMathOperator{\Span}{span}
\DeclareMathOperator{\Ker}{ker}
\DeclareMathOperator{\im}{im}
\DeclareMathOperator{\Hom}{Hom}
\DeclareMathOperator{\Aut}{Aut}
\DeclareMathOperator{\Var}{Var}
\DeclareMathOperator{\Cov}{Cov}
\newtheorem{problem}{Problem}[section]
\newenvironment{solution}{%
  \par\smallskip\noindent\textbf{Solution notes}\par
  \vspace{1.15in}\noindent\hrulefill\par\vspace{0.28in}%
}{\par}
\newenvironment{originalproblem}{%
  \VerbatimEnvironment
  \par\medskip\noindent\begin{minipage}{0.98\linewidth}
  \begin{Verbatim}[fontsize=\small]%
}{%
  \end{Verbatim}\end{minipage}\par
}
\renewcommand{\coursename}{Math 113}
\renewcommand{\assignment}{Homework 1 Answers}
\begin{document}
\begin{center}\Large\textbf{Math 113 -- Homework 1 Answers}\end{center}
\includepdf[pages=1,pagecommand={\thispagestyle{empty}}]{/Users/keshavramamurthy/Downloads/fa26_books/113_hw1.pdf}
\clearpage
\section{Expanded problems for credit}
\subsection*{Problem 1}
\begin{originalproblem}
1. Fraleigh Exercise 0.6. In Exercises 5 through 10, decide whether the object described is indeed a set (is well defined). Give an alternate description of each set:

{n \in \mathbb{Z} : n^2 < 0}.
\end{originalproblem}
\begin{solution}
% Write your proof or calculation here.
\end{solution}
\subsection*{Problem 2}
\begin{originalproblem}
2. Fraleigh Exercise 0.16. List the elements of the power set of the given set and give the cardinality of the power set:

(a) \emptyset, \quad (b) \{a\}, \quad (c) \{a,b\}, \quad (d) \{a,b,c\}.
\end{originalproblem}
\begin{solution}
% Write your proof or calculation here.
\end{solution}
\subsection*{Problem 3}
\begin{originalproblem}
3. Fraleigh Exercise 0.30. Determine whether the relation x R y on \mathbb{R}, defined by x R y if and only if x \ge y, is an equivalence relation. Describe the partition arising from the relation if it is an equivalence relation.
\end{originalproblem}
\begin{solution}
% Write your proof or calculation here.
\end{solution}
\subsection*{Problem 4}
\begin{originalproblem}
4. Fraleigh Exercise 2.2. Exercises 1 through 4 concern the binary operation * defined on S=\{a,b,c,d,e\} by means of Table 2.26. Compute (a*b)*c and a*(b*c). Can you say on the basis of this computations whether * is associative?
\end{originalproblem}
\begin{solution}
% Write your proof or calculation here.
\end{solution}
\subsection*{Problem 5}
\begin{originalproblem}
5. Fraleigh Exercise 2.8. Determine whether the binary operation * defined on \mathbb{Q} by a*b=ab+1 is commutative and whether it is associative.
\end{originalproblem}
\begin{solution}
% Write your proof or calculation here.
\end{solution}
\subsection*{Problem 6}
\begin{originalproblem}
6. Fraleigh Exercise 3.2. Determine whether the map \phi(n)=-n is an isomorphism from (\mathbb{Z},+) to (\mathbb{Z},+). If it is not an isomorphism, explain why not.
\end{originalproblem}
\begin{solution}
% Write your proof or calculation here.
\end{solution}
\subsection*{Problem 7}
\begin{originalproblem}
7. Fraleigh Exercise 3.30. In Exercises 29 through 32, give a careful proof for a skeptic that the indicated property of a binary structure \langle S,*\rangle is indeed a structural property. The operation * is associative.
\end{originalproblem}
\begin{solution}
% Write your proof or calculation here.
\end{solution}
\section{Practice problems}
\begin{originalproblem}
Practice problems (not for submission):
Fraleigh Exercises 0.5, 0.9, 0.11, 0.17, 0.25, 0.28, 0.29, 0.31, 0.36;
2.3, 2.17, 2.26, 2.28, 2.37; and 3.3, 3.4, 3.24, 3.25, 3.31, 3.33.
\end{originalproblem}
\end{document}
